\section*{Bloque A. Fundamentos teóricos y conceptuales}

\noindent\textit{Preguntas 1 a 30. Evalúan la comprensión analítica de los supuestos de la competencia perfecta, la maximización del beneficio, la oferta individual y de la industria, el equilibrio de corto y largo plazo, los excedentes y la eficiencia asignativa.}

\q{¿Cuál de las siguientes combinaciones describe mejor un mercado de competencia perfecta?}
\begin{opts}
  \item Pocas empresas, producto diferenciado y barreras a la entrada.
  \item Una sola empresa, producto único y poder de mercado.
  \item Muchos compradores y vendedores, producto homogéneo, información perfecta y libre entrada y salida.
  \item El Estado fija el precio y las cantidades de todas las empresas.
\end{opts}

\q{La curva de demanda que enfrenta una empresa perfectamente competitiva es:}
\begin{opts}
  \item perfectamente elástica al precio de mercado.
  \item descendente y con pendiente negativa.
  \item vertical.
  \item quebrada.
\end{opts}

\q{La curva de demanda del mercado en competencia perfecta es:}
\begin{opts}
  \item horizontal.
  \item vertical.
  \item descendente.
  \item perfectamente elástica.
\end{opts}

\q{La condición de maximización del beneficio de una empresa es:}
\begin{opts}
  \item $P=CMe$.
  \item $IMg=CMe$.
  \item $P=CVMe$.
  \item $IMg=CMg$.
\end{opts}

\q{Para una empresa perfectamente competitiva, el ingreso marginal es:}
\begin{opts}
  \item mayor que el precio.
  \item menor que el precio.
  \item igual al precio.
  \item igual a cero.
\end{opts}

\q{El beneficio económico de una empresa se define como:}
\begin{opts}
  \item el ingreso total menos los costos explícitos.
  \item el beneficio contable.
  \item el ingreso total menos los costos fijos.
  \item el ingreso total menos el costo económico total, que incluye costos explícitos e implícitos.
\end{opts}

\q{El beneficio normal se relaciona con el beneficio económico porque:}
\begin{opts}
  \item el beneficio económico siempre es positivo.
  \item el beneficio normal es el costo de oportunidad del capital; un beneficio económico nulo significa que se está obteniendo el beneficio normal.
  \item el beneficio normal coincide con el beneficio contable.
  \item el beneficio normal es un costo fijo irrelevante.
\end{opts}

\q{En el corto plazo, una empresa competitiva puede obtener:}
\begin{opts}
  \item siempre beneficios positivos.
  \item siempre beneficios nulos.
  \item siempre pérdidas.
  \item beneficios positivos, nulos o negativos, según la relación entre el precio y el costo medio.
\end{opts}

\q{En el corto plazo, la empresa competitiva suspende la producción cuando:}
\begin{opts}
  \item $P<CMe$.
  \item $P<CVMe$.
  \item $P<CMg$.
  \item $P<CFMe$.
\end{opts}

\q{La curva de oferta de una empresa competitiva en el corto plazo es:}
\begin{opts}
  \item la curva de costo medio.
  \item la curva de costo variable medio.
  \item la curva de demanda.
  \item el tramo de la curva de costo marginal situado por encima del mínimo del costo variable medio.
\end{opts}

\q{La curva de oferta de la industria competitiva se obtiene mediante:}
\begin{opts}
  \item la suma vertical de las curvas de oferta individuales.
  \item la suma horizontal de las curvas de oferta individuales.
  \item el promedio de las curvas de oferta individuales.
  \item el máximo de las curvas de oferta individuales.
\end{opts}

\q{En el equilibrio de largo plazo de una industria perfectamente competitiva con libre entrada y salida:}
\begin{opts}
  \item $P>CMe$ y hay beneficios extraordinarios.
  \item $P<CVMe$.
  \item $P=CMg=CMe$ y el beneficio económico es cero.
  \item $P=CMg>CMe$.
\end{opts}

\q{Si en una industria competitiva el precio es mayor que el costo medio mínimo, en el largo plazo:}
\begin{opts}
  \item las empresas salen y el precio sube.
  \item no hay entrada de empresas.
  \item el precio sube por falta de oferta.
  \item entran empresas, la oferta aumenta y el precio tiende a bajar.
\end{opts}

\q{Si en una industria competitiva el precio es menor que el costo medio mínimo, en el largo plazo:}
\begin{opts}
  \item salen empresas, la oferta disminuye y el precio tiende a subir.
  \item entran empresas y la oferta aumenta.
  \item el precio sigue bajando indefinidamente.
  \item no cambia nada.
\end{opts}

\q{En una industria de costo constante, la curva de oferta de largo plazo es:}
\begin{opts}
  \item ascendente.
  \item descendente.
  \item horizontal.
  \item vertical.
\end{opts}

\q{En una industria de costo creciente, la curva de oferta de largo plazo es:}
\begin{opts}
  \item ascendente.
  \item horizontal.
  \item descendente.
  \item vertical.
\end{opts}

\q{En una industria de costo decreciente, la curva de oferta de largo plazo es:}
\begin{opts}
  \item horizontal.
  \item ascendente.
  \item vertical.
  \item descendente.
\end{opts}

\q{El excedente del productor se define como:}
\begin{opts}
  \item el ingreso total menos el costo total.
  \item el beneficio contable.
  \item el ingreso total menos el costo variable.
  \item el excedente del consumidor.
\end{opts}

\q{El excedente del consumidor se define como:}
\begin{opts}
  \item la disposición a pagar menos el precio efectivamente pagado.
  \item el precio menos el costo marginal.
  \item el ingreso total menos el costo total.
  \item la recaudación impositiva.
\end{opts}

\q{El excedente total o bienestar social es:}
\begin{opts}
  \item la suma del excedente del consumidor y del productor, y se maximiza en el equilibrio competitivo.
  \item únicamente la recaudación del impuesto.
  \item la diferencia entre el excedente del productor y el del consumidor.
  \item cero en el equilibrio de mercado.
\end{opts}

\q{Un precio máximo fijado por debajo del precio de equilibrio genera:}
\begin{opts}
  \item un excedente de oferta.
  \item ningún efecto sobre el mercado.
  \item escasez y pérdida de bienestar.
  \item un aumento automático de la oferta.
\end{opts}

\q{Un precio mínimo fijado por encima del precio de equilibrio genera:}
\begin{opts}
  \item un excedente de oferta y pérdida de bienestar.
  \item escasez.
  \item ningún efecto.
  \item un aumento de la demanda.
\end{opts}

\q{La incidencia económica de un impuesto depende de:}
\begin{opts}
  \item las elasticidades relativas de oferta y demanda; el lado menos elástico soporta una mayor carga.
  \item exclusivamente de quién paga legalmente el impuesto.
  \item siempre de los consumidores.
  \item siempre de los productores.
\end{opts}

\q{En el caso de un subsidio, la incidencia económica implica que:}
\begin{opts}
  \item siempre se beneficia el consumidor.
  \item el lado menos elástico captura una mayor parte del subsidio.
  \item siempre se beneficia el productor.
  \item no hay efectos sobre el precio.
\end{opts}

\q{La eficiencia asignativa en competencia perfecta se alcanza cuando:}
\begin{opts}
  \item $P=CMe$.
  \item $P=CVMe$.
  \item $P=CMg$.
  \item $P=0$.
\end{opts}

\q{La eficiencia productiva en el largo plazo competitivo se alcanza cuando:}
\begin{opts}
  \item $P>CMe$.
  \item $P=CMg>CMe$.
  \item $P=\min CMe$.
  \item $P=CVMe$.
\end{opts}

\q{La libre entrada de empresas en una industria competitiva:}
\begin{opts}
  \item aumenta los beneficios extraordinarios de las empresas existentes.
  \item no tiene efecto sobre los beneficios.
  \item crea un monopolio.
  \item elimina los beneficios económicos en el largo plazo.
\end{opts}

\q{Una empresa es precio-aceptante cuando:}
\begin{opts}
  \item fija un precio superior al de mercado.
  \item enfrenta una demanda descendente.
  \item puede influir en el precio mediante su nivel de producción.
  \item no puede influir en el precio de mercado y lo toma como dado.
\end{opts}

\q{El equilibrio de mercado en competencia perfecta se determina cuando:}
\begin{opts}
  \item $Q_d=Q_s$.
  \item $Q_d>Q_s$.
  \item $Q_d<Q_s$.
  \item $P=0$.
\end{opts}

\q{En el corto plazo, una empresa competitiva produce:}
\begin{opts}
  \item donde $P=CMg$, siempre que $P\geq CVMe$.
  \item siempre en el punto de cierre.
  \item donde $P=CMe$.
  \item fijando un precio superior al de mercado.
\end{opts}
